% ============================================================================
%  Actividad 4 · Alcance declarado del prototipo
%
%  Propiedad de US-UX-07. Responde al criterio 5 de la historia: la tabla de
%  alcance de tres estados que sustituye a la tabla de paridad original.
%
%  Regla que gobierna este archivo: la primera tabla publica el valor que el
%  codigo declara en PROTOTIPOS, no uno deseado. frontend/test/alcancePrototipos.spec.ts
%  compara ambas cosas y falla si divergen.
% ============================================================================

% ============================================================================
\uxsection{Alcance declarado de cada pantalla}

Un prototipo de alta fidelidad se parece a un producto terminado, y ese parecido es justamente lo
que obliga a declarar hasta dónde llega. Esta sección dice, pantalla por pantalla y actividad por
actividad, qué está construido, qué está construido sin datos detrás y qué está solamente escrito.
El criterio es el mismo que la Actividad 3 aplicó a su wireframe: no presentar diseño como producto
en funcionamiento.

El esquema tiene tres estados y solo tres:

\begin{uxlist}
  \lead{Navegable con datos de ejemplo}{La pantalla se recorre y muestra contenido proveniente de
    los silos sintéticos o de la base del portal. Lo que se ve es el resultado de una consulta real
    sobre datos reproducibles, no un texto tecleado en la plantilla.}
  \lead{Navegable sin datos}{La pantalla se recorre, su composición es la definitiva y sus estados
    están resueltos, pero el contenido que llenaría la caja todavía no proviene de una fuente. La
    pantalla lo declara en su propio cuerpo y nombra la historia de usuario que lo entrega.}
  \lead{En hoja de ruta}{La capacidad está descrita y no está construida. Aparece en este documento para
    que nadie tenga que deducir su ausencia.}
\end{uxlist}

\subsection{Las siete pantallas}

La tabla siguiente publica, para cada prototipo, el valor que el contrato de navegación declara en
\texttt{frontend/app/utils/navegacion.ts}. No es una lectura de quien redacta: es el literal que la
aplicación usa para pintar la insignia de alcance en el índice, y una prueba automática compara
esta tabla contra ese literal en cada corrida.

% tabla-alcance:inicio
% Las siete filas de abajo las lee frontend/test/alcancePrototipos.spec.ts y las
% compara contra PROTOTIPOS. El formato es contrato: la primera celda es la ruta
% dentro de \texttt{} y la tercera es la etiqueta de alcance. No reordenar las
% columnas sin ajustar la prueba.
\begin{uxtabla}{|L{3.5cm}|L{2.6cm}|L{3.1cm}|Y|}{Alcance declarado de los siete prototipos}
  \uxheadrow \thd{Ruta} & \thd{Pantalla} & \thd{Alcance} & \thd{Qué falta para el estado siguiente} \\
  \texttt{/acceso} & Acceso & Navegable sin datos & El formulario de usuario y contraseña opera contra el emisor de credenciales; la puerta de demostración es la vía de la evaluación \\
  \texttt{/inicio} & Inicio & Navegable sin datos & El buscador, las búsquedas recientes, las favoritas y las alertas componen desde marcadores rotulados como datos de ejemplo, no desde el catálogo \\
  \texttt{/exploracion} & Exploración y extracción & Navegable sin datos & El catálogo temático y los filtros facetados los entrega US-008; la pantalla declara las cuatro capacidades y la historia que las trae \\
  \texttt{/gobierno} & Gobierno del dato & Navegable sin datos & El diccionario de campos y el linaje se recorren; la procedencia todavía no se resuelve contra la base de catálogo \\
  \texttt{/asistente} & Asistente conversacional & Navegable sin datos & El transporte, las tarjetas de llamada a herramienta y la cancelación son reales; el contenido proviene de un proveedor guionizado \\
  \texttt{/administracion} & Administración & Navegable sin datos & El panel de personas usuarias se recorre; la edición completa queda degradada por acuerdo del sprint \\
  \texttt{/exploracion/exportar} & Exportación & Navegable sin datos & El ciclo de vida del trabajo de exportación es real contra el api; el almacenamiento definitivo del archivo queda fuera de esta semana \\
\end{uxtabla}
% tabla-alcance:fin

\subsection{Una divergencia declarada, no corregida}

Las siete filas dicen \emph{navegable sin datos} y varias pantallas muestran más de lo que esa
etiqueta promete: la de inicio compone cinco bloques con contenido visible, la de gobierno recorre
un diccionario de campos y la de administración pinta un panel de personas usuarias. La etiqueta
del código quedó fijada antes de que esas historias cerraran y nadie la reescribió durante el
sprint, porque el contrato de navegación es un archivo que varias historias leen y ninguna tiene
asignado en escritura esta semana.

Se publica el valor del código y no el observado. El motivo es que esta tabla vale exactamente por
su verificabilidad: si el documento calificado y la aplicación pueden divergir por una edición de
última hora que nadie más está esperando, la tabla deja de ser evidencia y pasa a ser una opinión
sobre el producto. La corrección de las etiquetas queda anotada como pendiente para la Actividad 5,
donde el contrato de navegación sí tiene una historia dueña.

% ============================================================================
\section{Traducción del estado de las historias de usuario}

La tabla anterior habla de pantallas y el equipo lleva su registro por historias de usuario, con un
vocabulario más fino. Esta segunda tabla traduce ese registro al esquema de tres estados, de modo
que ninguna historia quede sin declarar. La correspondencia es fija:

\begin{uxtabla}{|L{6.2cm}|Y|}{Correspondencia entre el registro del equipo y el esquema publicado}
  \uxheadrow \thd{Estado en el registro del sprint} & \thd{Estado publicado} \\
  Cerrada & Navegable con datos de ejemplo \\
  Cerrada degradada & Navegable con datos de ejemplo \\
  Demostrada en prototipo, no productiva & Navegable sin datos \\
  Roadmap declarado en el documento & En hoja de ruta \\
  Congelada, sin cambio & En hoja de ruta \\
  De la Actividad 5, no se toca & En hoja de ruta \\
\end{uxtabla}

El registro del sprint enuncia cuatro estados y enumera seis filas. Se publican las seis, que son
las que el registro usa en la práctica; la discrepancia entre el enunciado y la enumeración queda
anotada aquí y se resuelve en el registro, no en este documento.

\begin{uxtablalarga}{|L{2.6cm}|L{3.6cm}|Y|}{Estado declarado de todas las historias de usuario al cierre del sprint}{\thd{Historia} & \thd{Estado publicado} & \thd{Precisión}}
  US-001 & Navegable con datos de ejemplo & Entorno de desarrollo y contrato de navegación \\
  US-002 & Navegable con datos de ejemplo & Cerrada sin integración continua \\
  US-003 & Navegable con datos de ejemplo & Puente de despliegue, sin infraestructura como código \\
  US-006 & Navegable con datos de ejemplo & Volúmenes de datos recortados respecto de lo planeado \\
  US-008 & Navegable con datos de ejemplo & Búsqueda por palabra clave únicamente \\
  US-009 & Navegable sin datos & Exportación demostrada en prototipo, no productiva \\
  US-011 & En hoja de ruta & La versión mínima alimenta las pantallas; las diez consultas de referencia quedan fuera \\
  US-012, US-013, US-014 & En hoja de ruta & Congeladas sin cambio \\
  US-015 & Navegable con datos de ejemplo & Autenticación y credenciales \\
  US-016, US-017 & Navegable con datos de ejemplo & Autorización por rol y visibilidad de módulos \\
  US-018, US-019 & Navegable con datos de ejemplo & Administración de personas usuarias sin edición completa \\
  US-020, US-021 & En hoja de ruta & Agente conversacional con modelo real. El proveedor del portal es guionizado \\
  US-022 & En hoja de ruta & Congelada sin cambio \\
  US-023 & Navegable sin datos & Transmisión y cancelación reales; contenido guionizado \\
  US-024 & Navegable con datos de ejemplo & Error tipado en la transmisión, sin control de reintento \\
  US-025 & Navegable con datos de ejemplo & Serie de alto volumen, con el tope acordado en el sprint \\
  US-026, US-027 & Navegable con datos de ejemplo & Linaje y bitácora \\
  US-028 & Navegable con datos de ejemplo & Tarjetas de visibilidad de llamada a herramienta \\
  US-029 & Navegable con datos de ejemplo & Estado compartido entre tablero y asistente, en versión estática \\
  US-030 a US-034 & En hoja de ruta & Pertenecen a la Actividad 5: observabilidad, tiempo al primer token, costo y paso a producción \\
  US-035, US-036 & En hoja de ruta & Congeladas sin cambio \\
  US-037 a US-041 & En hoja de ruta & Declaradas fuera del alcance del prototipo \\
  US-004, US-005 & En hoja de ruta & Pertenecen a la Actividad 5 \\
  US-UX-07, US-UX-09 & Navegable con datos de ejemplo & Esta actividad y su guía de estilos \\
  US-UX-08 & En hoja de ruta & Pertenece a la Actividad 5 \\
  Bloques de holgura & En hoja de ruta & Congelados desde la planeación \\
\end{uxtablalarga}

% ============================================================================
\section{Las seis capacidades que la Actividad 2 prometió}

El journey map de la Actividad 2 identificó oportunidades de diseño que el catálogo de historias no
llegó a cubrir. Un evaluador puede contrastar aquellas promesas contra este prototipo, así que se
declaran una por una. Ninguna está construida.

\begin{uxtabla}{|L{5.0cm}|L{2.2cm}|Y|}{Capacidades prometidas en la Actividad 2 y su estado}
  \uxheadrow \thd{Capacidad} & \thd{Estado} & \thd{Dónde aparecería} \\
  Consultas guardadas & En hoja de ruta & Exploración y extracción, junto a la consulta con filtros \\
  Marca de versión retirada con enlace a la vigente & En hoja de ruta & Gobierno del dato, como atributo del campo en el diccionario \\
  Notificación dirigida por cambio de definición & En hoja de ruta & Inicio, en el bloque de alertas \\
  Copiado con procedencia & En hoja de ruta & Transversal: acompañaría a toda cifra copiada desde una tabla \\
  Comparación lado a lado de variables similares & En hoja de ruta & Exploración y extracción \\
  Credenciales de interfaz de programación autogestionadas & En hoja de ruta & Administración, como faceta transversal \\
\end{uxtabla}

Dos de las seis son baratas y de alto peso narrativo, porque se apoyan en atributos que el catálogo
ya modela: la marca de versión retirada y el copiado con procedencia. Quedan propuestas para la
Actividad 5 con esa justificación. Las cuatro restantes exigen trabajo de producto que ninguna
historia del catálogo cubre hoy.

El silencio sobre estas seis capacidades sería el único desenlace sin recurso: quien lea la
Actividad 2 y recorra este prototipo encontraría la diferencia por su cuenta y sin explicación. Por
eso se declaran aquí, con nombre y con estado.
